A computer is a machine for discrete mathematics. Every program is a formula of logic, every data structure a graph,
every search a walk through an enormous finite space. This book is a journey through the mathematics behind them,
told through the problems and the people who solved them.

Why must four of any five points form a convex quadrilateral? How many isomers does a molecule of thirty carbon
atoms have? Why can a map always be coloured with five colours, and when can every applicant find a job? How large
can a set be without containing three numbers in arithmetic progression? Along the way: logic and SAT solving,
orders and fixed points, counting with symmetry, generating functions, graphs and networks, matchings and flows,
planarity, Ramsey theory, eigenvalues of graphs, matroids, the linear algebra method and finite geometries.

Every theorem is proved in full, every exercise has a hint and a reasoned solution, and every chapter comes with a
Python notebook that turns the theorems into working algorithms. Written for the pleasure of seeing why things are
true.
